\begin{lemma}
\label{lemma-difficulty-K-injectives}
Let $\mathcal{A}$ be an abelian category with countable products and
enough injectives. Let $K^\bullet$ be a complex. Let $I_n^\bullet$ be
the inverse system of bounded below complexes of injectives produced by
Lemma \ref{lemma-special-inverse-system}. Then
$I^\bullet = \lim I_n^\bullet$ exists, is K-injective, represents
$R\lim \tau_{\geq -n}K^\bullet$ in $D(\mathcal{A})$,
and the following are equivalent
\begin{enumerate}
\item the map $K^\bullet \to I^\bullet$ (see proof) is a quasi-isomorphism,
\item the map $K^\bullet \to R\lim \tau_{\geq -n}K^\bullet$
of Remark \ref{remark-map-into-derived-limit-truncations}
is an isomorphism in $D(\mathcal{A})$.
\end{enumerate}
\end{lemma}

\begin{proof}
The statement of the lemma makes sense as $R\lim \tau_{\geq -n}K^\bullet$
exists by Lemma \ref{lemma-inverse-limit-bounded-below}.
Each complex $I_n^\bullet$ is K-injective by
Lemma \ref{lemma-bounded-below-injectives-K-injective}.
Choose direct sum decompositions $I_{n + 1}^p = C_{n + 1}^p \oplus I_n^p$
for all $n \geq 1$. Set $C_1^p = I_1^p$. The complex
$I^\bullet = \lim I_n^\bullet$ exists because we can take
$I^p = \prod_{n \geq 1} C_n^p$. Fix $p \in \mathbf{Z}$.
We claim there is a split short exact sequence
$$
0 \to I^p \to \prod I_n^p \to \prod I_n^p \to 0
$$
of objects of $\mathcal{A}$. Here the first map is given by
the projection maps $I^p \to I_n^p$ and the second map
by $(x_n) \mapsto (x_n - f^p_{n + 1}(x_{n + 1}))$ where
$f^p_n : I_n^p \to I_{n - 1}^p$ are the transition maps.
The splitting comes from the map $\prod I_n^p \to \prod C_n^p = I^p$.
We obtain a termwise split short exact sequence of complexes
$$
0 \to I^\bullet \to \prod I_n^\bullet \to \prod I_n^\bullet \to 0
$$
Hence a corresponding distinguished triangle in $K(\mathcal{A})$
and $D(\mathcal{A})$. By Lemma \ref{lemma-product-K-injective}
the products are K-injective and represent the corresponding
products in $D(\mathcal{A})$.
It follows that $I^\bullet$ represents $R\lim I_n^\bullet$
(Definition \ref{definition-derived-limit}).
Since $R\lim I_n^\bullet \cong R\lim \tau_{\geq -n}K^\bullet$
as derived limits are defined on the level of the derived
category, we see that $I^\bullet$ represents $R\lim \tau_{\geq -n}K^\bullet$.
Moreover, the complex $I^\bullet$ is K-injective by
Lemma \ref{lemma-triangle-K-injective}.
By the commutative diagram of Lemma \ref{lemma-special-inverse-system}
and since $K^i = (\tau_{\geq -n}K^\bullet)^i$
for $n \gg 0$ we see that we get a unique map
$\gamma : K^\bullet \to I^\bullet$ such that the diagrams
$$
\xymatrix{
K^\bullet \ar[r] \ar[d]_\gamma & \tau_{\geq -n} K^\bullet \ar[d] \\
I^\bullet \ar[r] & I_n^\bullet
}
$$
commute. It follows that $\gamma$ is a map of complexes which represents
the map $c : K^\bullet \to R\lim \tau_{\geq -n}K^\bullet$ of
Remark \ref{remark-map-into-derived-limit-truncations}
in $D(\mathcal{A})$. In other words, the diagram
$$
\xymatrix{
K^\bullet \ar[r]_-c \ar[d]_\gamma & R\lim \tau_{\geq -n} K^\bullet
\ar[d]^{\cong} \\
I^\bullet \ar[r]^-{\cong} & R\lim I_n^\bullet
}
$$
is commutative in $D(\mathcal{A})$. The lemma follows.
\end{proof}